\section*{Что считается атомистическим равновесием}
В расчёте использован тот же опубликованный потенциал моментных тензоров
Qi и соавт., что и в серии при 1000~К. Состояние содержит типы всех
узлов, координаты атомов, объём и отношение осей гексагональной ячейки.
При заданном \(T,\ p=0\) в ходе Монте-Карло предлагаются замены
Ti\(\rightleftarrows\)Al, смещения атомов, изменения объёма и формы.
Химические степени свободы и ангармонические колебания входят в один
классический ансамбль. Для смены объёма учитывается множитель
\((V'/V)^N\). Один проход содержит \(N\) попыток химической замены и
\(N\) попыток смещения. Объём ограничен диапазоном 15--21~\AA$^3$/атом,
\(c/a\) в основном расчёте меняется в диапазоне 1,45--1,72.
В механическом контроле фиксируется авторское \(c/a=1{,}56374\);
объём продолжает флуктуировать. Каждой изотерме соответствуют ячейки
\(4\times4\times3\), \(6\times6\times4\), \(8\times8\times6\) и
\(10\times10\times7\) из двухатомных ячеек ГПУ.
\href{https://doi.org/10.1103/PhysRevMaterials.7.103602}{Qi et al. (2023).}

Исследуется равновесие кристаллических \(\alpha\) и \(\alpha_2\) на
ГПУ-решётке. Чтобы исключить ранее обнаруженное нефизическое сближение
Ti--Ti у исходного MTP, для каждого атома ограничена проекция смещения
от идеального узла вдоль каждой из 12 ближайших связей ГПУ. Величина
\(W\) --- максимум этих проекций, нормированных на половину длины соответствующей связи,
после удаления общего переноса. Основному расчёту соответствует
\(W<0{,}40\). Для объёмных и осевых проб дополнительно требуется
\(d_{\rm nn}^{\rm ideal}(V,c/a)(1-W_{\max})\ge1{,}68052\)~\AA,
где \(d_{\rm nn}^{\rm ideal}\) --- кратчайшее расстояние между идеальными узлами.
В рамках такого определения бассейна возможная конкуренция
со структурами вне ГПУ остаётся отдельной задачей.

Пусть \(n_j\) --- число атомов Al на каждой из четырёх подрешёток,
\(k=\sum_jn_j\). Параметр дальнего порядка
\[
 Q=\frac{4(4\max_j n_j-k)}{3N}
\]
равен единице для идеального Ti$_3$Al типа D0$_{19}$. В конечной
разупорядоченной ячейке \(Q\) положителен из-за операции максимума.
В окне \(l\) к энергии прибавлялся член
\(\frac12 K_l(Q-q_l)^2-\Delta\mu_l k\).
Химическое поле \(\Delta\mu_l\) относится к конфигурационной замене;
пересчёт к физическим массам даёт общий сдвиг поля при данной
температуре и сохраняет равновесные составы.

По измеренным \(k,Q\) метод MBAR снимает зонтики и восстанавливает
относительные веса состояний. При одинаковых \(N,T,W,p\), механическом
условии на \(c/a\) и допустимом пространстве конфигураций общая
атомистическая энергия сокращается в отношениях весов между окнами.
Это определение относительных свободных энергий позволяет обойтись
без абсолютной свободной энергии референсного кристалла. Для геометрически
допустимых проб изменение энергии вычисляется непосредственно по MTP. В каждой
изотерме используется собственное \(k_{\rm B}T\).
Когда две моды разделены минимумом \(Q^\ddagger\), порядок складывает
четыре эквивалентных варианта D0$_{19}$. В конечной ячейке оценка поля
сосуществования для одного варианта упорядочения определяется из уравнения
\[
 P_{\Delta\mu^*}(Q>Q^\ddagger)=4P_{\Delta\mu^*}(Q\le Q^\ddagger).
\]
По обе стороны разделителя вычисляют условные средние
\(\langle k/N\rangle\). Блочное ресэмплирование сохраняет временные
корреляции внутри траекторий. Получаемые интервалы условны по отношению
к размеру ячейки, определению бассейна и покрытию ансамбля окнами.
Если два отдельных бассейна не разрешены, два произвольных среза по \(Q\)
дали бы разные составы даже у одной фазы. В этой ситуации анализируют
изотерму состава, моды \(P(Q)\), восприимчивость и зависимость от размера.
\href{https://doi.org/10.1063/1.2978177}{Shirts, Chodera (2008).}

\subsection*{Граница применимости эталона}
На 288 узлах при 1273,15~К и поле 3,48~эВ уменьшение \(W_{\max}\)
с 0,40 до 0,35 повысило \(\langle Q\rangle\) с 0,29 до 0,40,
а среднюю концентрацию --- на 0,31 п.п. Таким образом, определение
кристаллического бассейна заметно для степени порядка. Соответствующие
профили и расчёт при фиксированном \(c/a\) проверены отдельно. Перенос
результата на неограниченный потенциал требует устранения нефизической
короткодействующей ветви и повторной проверки свободных энергий.

